\documentclass[12pt]{article}
\usepackage{latexsym, amscd, %graphicx, 
amssymb}
\newcommand{\SECT}[1]
		{\vspace{.3cm}
		$\!\!\!\!\!\!\!\!\!\!\!${\bf {#1}}
		\vspace{.2cm}}

\begin{document}


\begin{center}

{\large \bf Math 152, Fall, 2004, Workshop 5\par}
\vspace{.5cm}
{{\bf Honors Section} 
\par}
\end{center}



\begin{enumerate}


\item \begin{enumerate}

\item Use the Trapezoidal Rule with $n=4$ to approximate 
$$\int_{0}^{1}-\cos(x^{2}+3) dx.$$
Your answer 
should contain $4$ digits past the decimal point (e.g. $1.2345$).

\item How good is the approximation obtained above?

\item How large must $n$ be so that the approximation given by the 
Trapezoidal Rule is correct to $4$ digits past the decimal point 
(i.e. so that $|E_{T_{n}}|<0.0001$)? 

\end{enumerate}

\noindent ({\bf Some useful inequalities:} 

\begin{enumerate}

\item {\bf Triangle inequality} $|f(x)+g(x)|\le |f(x)|+|g(x)|$; 
$|f(x)-g(x)|\le |f(x)|+|g(x)|$.

\item If $|f(x)|\le F$ and $H\le |g(x)|\le G$, then $|f(x)g(x)|\le FG$
and ${\displaystyle \frac{|f(x)|}{|g(x)|}\le \frac{F}{H}}$.

\item If $f(x)$ is increasing on $[a, b]$, then $f(x)\le f(b)$;
if $f(x)$ is decreasing on $[a, b]$, then $f(x)\le f(a)$.

\item $|\sin x|\le 1$ and $|\cos x|\le 1$ for all real numbers.

\end{enumerate}
These principles can be used to find an upper bound for a function quickly,
especially when trying use the Differential Calculus to maximize that 
function is likely to be quite arduous. )


\item Assume that $a$ is a positive constant. Let $R$ be the region 
bounded above by $y = 1/x^{a}$, below 
by $y = 0$, and on the left by the line $x = 1$. 
Determine those values of $a$ for which the integral 
that results from attempting to calculate each of the following converges: 

\begin{enumerate}

\item The volume of the solid obtained by rotating $R$ around the $x$-axis.

\item The volume of the solid obtained by rotating $R$ around the $y$-axis. 


\end{enumerate}

\item  Suppose that the functions $f$ and $g$ are defined and 
positive on some interval $[a, \infty)$. 

\begin{itemize}

\item If ${\displaystyle \lim_{x\to \infty} \frac{f(x)}{g(x)}
= 0}$, we say that $f$ grows more slowly than $g$ or that $g$ 
grows more rapidly than 
$f$ , as $x \to \infty$. This may be indicated by writing 
$f << g$. (We use this terminology even if $f$ 
and/or $g$ do not in fact ``grow" as $x \to \infty$; for example, 
they might go to zero). 

\end{itemize}

\begin{enumerate}

\item Arrange the following eight functions in order of rate 
of growth, from least rapidly growing 
to most rapidly growing. indicating any cases in which the 
functions grow at comparable rates. 
Justify your answers. It is helpful to least check that if 
$f << g$ and $g << h$ then $f << h$. 
$$x; \ln x;  \sqrt{x^{3} + x}; e^{x} ; \ln(\ln x); \ln (e^{x^{2}}
+ 1); x \ln(x); e^{x^{2}}.$$

\item Arrange the reciprocals of the same eight functions in order 
of rate of growth. (Hint: if $f$ and $g$ 
are related by $f << g$, what is the relation between $1/g$ and $1/f$?)
 
\item For which of the reciprocals $h(x) = 1/f(x)$ in (b) does the integral 
${\displaystyle \int_{100}^{\infty}
h(x) dx}$ converge? (Suggestion: Use the Comparison Theorem 
on page 536 of the text. To do so you will need to 
answer this queston: if $f << g$ is it true that $f(x) \le g(x)$? 
Or true for some $x$?)

\end{enumerate}






\end{enumerate}



\end{document}


